
%%%
\subsection{Vectors}

\subsubsection{\definition{Linear Combination} $\lambda \mathbf{v} + \mu\mathbf{w}$}
$\mathbf{v},\mathbf{w} \in \mathbb{R}^m$, $\lambda, \mu \in \mathbb{R}$

$\sum^n_{i=1}\lambda_i\mathbf{v}_i$, e.g. $5\begin{pmatrix}2\\3\end{pmatrix} - 3\begin{pmatrix}3\\-1\end{pmatrix} =
\begin{pmatrix}10\\15\end{pmatrix} - \begin{pmatrix}9\\-3\end{pmatrix} =
\begin{pmatrix}1\\18\end{pmatrix}$

Every vector in $\mathbb{R}^m$ can be written as $\sum^m_{i=1} u_i\mathbf{e}_i$.

\begin{enumerate}[label=(\roman*)]
\item 
    \textbf{affine}: if $\lambda_1 + \lambda_2 + ... + \lambda_n = 1$
\item 
    \textbf{conic}: if $\lambda_j \geq 0$ for $j=1,2,...,n$
\item 
    \textbf{convex}: if both \textit{affine} and \textit{conic} combination
\end{enumerate}
\subsubsection{Lengths and Angles from Dot Products}
\definition{Scalar (or dot, inner) product}:
$\mathbf{v} \cdot \mathbf{w} = v_1w_1 + \cdots + v_nw_n$. \\
\definition{Outer product} $rank(A)=1 \iff \exists \text{ non-zero vectors }\mathbf{v} \in \mathbb{R}^m, \mathbf{w} \in \mathbb{R}^n$, s.t. $A$ is an outer product, i.e. $A=\mathbf{vw}^\top$, thus $rank(\mathbf{vw}^\top)=1$.
\\[2.22mm]
\definition{Length} (Euclidian Norm): $||\mathbf{v}|| = \sqrt{\mathbf{v} \cdot \mathbf{v}} = \sqrt{\mathbf{v} ^\top \mathbf{v}} = \sqrt{v^2_1 + \cdots + v^2_n}$. $||\mathbf{v}||^2:=\mathbf{v}^\top \mathbf{v}$\\
\textbf{\definition{Unit vector}}: $||\mathbf{u}||=1=\frac{\mathbf{v}}{||\mathbf{v}||}$, for $\mathbf{v} \ne \mathbf{0}$.
\\[2.22mm]
\definition{Orthogonal} (perpendicular) vectors: $\mathbf{v} \cdot \mathbf{w} = 0$. \\
\definition{Angle between vectors} (Cosine Formula): $\cos{\alpha} = \frac{\mathbf{v} \cdot \mathbf{w}}{||\mathbf{v}||||\mathbf{w}||}$
, for $\mathbf{v} \ne \mathbf{0}$. \\
Because $\cos{\alpha} \le 1$, \lemma{Cauchy-Schwarz} inequality:
$$\underbrace{|\mathbf{v} \cdot \mathbf{w}|}_{|\cos{\alpha}|||\mathbf{v}||||\mathbf{w}||}
\le ||\mathbf{v}||||\mathbf{w}||.$$
\lemma{Triangle inequality}: $||\mathbf{v} + \mathbf{w}|| \le ||\mathbf{v}|| + ||\mathbf{w}||$.

\subsubsection{\definition{Linear (In)dependence} of Vectors}
$\mathbf{w}_1, \mathbf{w}_2, \dots , \mathbf{w}_k$ are linearly \textbf{independent} if:
\begin{enumerate}[label=(\roman*)]
\item 
    \underline{no} vector is a linear combination of the previous ones. Or
\item 
    \underline{no} vector is a linear combination of the other ones. Or
\item 
    there are \underline{no} $ c_1, c_2, \dots , c_k$ besides $0, 0, \dots , 0$ such that
    $c_1\mathbf{w}_1 + c_2\mathbf{w}_2 + \dots + c_k\mathbf{w}_k = \mathbf{0}.$
\end{enumerate}
All are equivalent: $(i) \Leftrightarrow (ii) \Leftrightarrow \ (iii)$. \\
Replace no by \textbf{\underline{some}} for linear \textbf{dependence}.
\\[2.22mm]
$A$ has \textbf{independent columns} if (iii) there is no $\mathbf{x}$ besides $\mathbf{0}$ such that $A\mathbf{x} = \mathbf{0}$.

\subsubsection{\definition{Span}}
\textbf{Span} of vectors is the set of all linear combinations of those vectors.


\subsubsection{\definition{Hyperplane through origin}}
Let $\mathbf{d}\in \mathbb{R}^m$, $\mathbf{d} \ne 0$, $H_{\mathbf{d}} = \{\mathbf{v} \in \mathbb{R}^m : \mathbf{v} \cdot \mathbf{d} = 0\}$

%%%
\subsection{Matrices Basics}

$A \in \mathbb{R}^{m \times n}$: a \definition{matrix} $A$ with $m$ rows, $n$ columns. {\tiny \textit{\textbf{Z}eilen \textbf{z}uerst, \textbf{S}palten \textbf{s}päter}}
$A=[a_{ij}]^{m\ \ \ \ \ n}_{i=1,\ j=1}$\\
\definition{\textbf{Trace}}: Sum of the diagonal entries.\\
$\definition{\textbf{rank}(A)} =$ \# independent columns (rows).


%%%

\subsubsection{Distributivity and associativity}
$A(B+C) = AB+AC$, $(B+C)D = BD+CD$.
And $(AB)C = A(BC) = ABC$, $(AB)(CD) = A((BC)D) = \cdots = ABCD$.
Matrix multiplication isn't commutative.


%%%
\subsection{The \definition{Transpose} of $A$}
%\subsubsection{The Transpose of $A$}
\begin{multicols}{2}
$$
\begin{bmatrix}
    a & b \\
    c & d \\
    e & f \\
\end{bmatrix}^\top
=
\begin{bmatrix}
    a & c & e\\
    b & d & f\\
\end{bmatrix}
$$
Row $i$ of $A =$ column $i$ of $A^\top$, column $j$ of $A =$ row $j$ of $A^\top$.
\end{multicols} 
\noindent\textbf{Product:} $(AB)^\top = B^\top A^\top$. \\
More matrices: $(ABC)^\top = C^\top B^\top A^\top$. \\[2.22mm]
\textbf{Inverse:} $(A^{-1})^\top = (A^\top)^{-1}$. \\
\textit{Proof:} $(A^{-1})^\top A^\top = (AA^{-1})^\top = I^\top = I$. \\[2.22mm]
\textbf{Permutation}: $P^\top = P^{-1}$. \\
\textit{Proof:} $\mathbf{p}_i \cdot \mathbf{p}_j = I_{ij} \Leftrightarrow PP^\top = I$.

%%%
\subsection{{Inverse Matrices}}
$M \in \mathbb{R}^{n \times n}$ (square!) is invertible if there is a matrix $M^{-1} \in \mathbb{R}^{n \times n}$ (the inverse of $M$) such
that $MM^{-1} = M^{-1}M = I$. \\
There can only be \textbf{one inverse}: If $MX = YM = I$, then $X = IX = (YM)X = Y(MX) = YI = Y$. \\[2.22mm]
$\mathbb{R}^{1 \times 1}$: $\begin{bmatrix}x\end{bmatrix}^{-1} = \begin{bmatrix}\frac{1}{x}\end{bmatrix}$ (if $x \not= 0$) \\
$\mathbb{R}^{2 \times 2}$: $
\begin{bmatrix}
    a & b \\
    c & d \\
\end{bmatrix}^{-1} = \displaystyle\frac{1}{ad-bc}
\begin{bmatrix}
    d & -b \\
    -c & a \\
\end{bmatrix}$ if ($ad-bc \not= 0$) \\
$\mathbb{R}^{n \times n}$: $[M|I]\rightarrow[I|M^{-1}]$ using Gauss-Jordan elimination (3.2.1)\\

\subsubsection{The Inverse Theorem}
$\phantom{\Leftrightarrow{ }}A \in \mathbb{R}^{n\times n}$ is invertible \\
$\Leftrightarrow \forall b \in \mathbb{R}^n$, $A\mathbf{x} = \mathbf{0}$ has a unique solution $\mathbf{x}$ \\
$\Leftrightarrow$ The columns of A are independent. \\[2.22mm]
For any $A \in \mathbb{R}^{n\times n}, B \in \mathbb{R}^{n\times n}$, if $AB = I$, then $BA = I$. \\
\textit{Proof:} $AB = I \Leftrightarrow B = A^{-1} \Leftrightarrow BA = A^{-1}A = I$.

%% TODO: Determinant

\subsubsection{Properties of the inverse}
If $A \in \mathbb{R}^{n\times n}$ and $B \in \mathbb{R}^{n\times n}$ are invertible, then
$(AB)^{-1} = B^{-1}A^{-1}$.
\textit{Proof:} $(AB)(B^{-1}A^{-1}) = A(BB^{-1})A^{-1} = AIA^{-1} = I$. \\
More matrices: $(ABC)^{-1} = C^{-1}B^{-1}A^{-1}$. \\
$(A^{-1})^{-1}=A$ \\
$(A^\top)^{-1}=(A^{-1})^\top$

\subsection{\textcolor{red}{Special Matrices}}
\begin{itemize}
    \item \textbf{Identity Matrix} $(a_{ii}=1\text{ for all }i)$: $I$
    \item \textbf{Diagonal Matrix} $(a_{ij}=0\text{ for all }i\ne j)$
    \item \textbf{Upper triangular matrix} $(a_{ij}=0\text{ for all }i>j):U$
    \item \textbf{Lower triangular matrix} $(a_{ij}=0\text{ for all }i<j):L$
    \item \textbf{Symmetric matrix} $(a_{ij}=a_{ji}\text{ for all }i,j):A=A^\top$
    \item \textbf{Skew-symmetric matrix} $(a_{ij}=-a_{ji}\text{ for all }i,j):A=-A^\top$
    \item \textbf{Nilpotent matrix} $A^k = 0\text{ for some }k \in \mathbb{N}$
\end{itemize}

\subsubsection{Symmetric Matrices}
A matrix $S_{n\times n}$ satisfying $S = S^\top$ is symmetric. \\[2.22mm]
If $S$ is symmetric, then $S^{-1}$ is also symmetric (if it exists). \\
\textit{Proof:} $S^{-1} = (S^\top)^{-1} = (S^{-1})^\top$. \\[2.22mm]
$AA^\top$ (and $A^\top A$) is symmetric. \\
\textit{Proof:} $AA^\top = (A^\top)^\top A^\top = (AA^\top)^\top$.

%%
\subsection{\theorem{$CR$-Decomposition} $A=CR$}
$C \in \mathbb{R}^{m \times r}$: the independent columns. \\
$R \in \mathbb{R}^{r \times n}$: how to combine them to get all columns. \\
$r = rank(A)$
$$
\begin{bmatrix}
    1 & 2 & 0 & 3 \\
    2 & 4 & 1 & 4 \\
    3 & 6 & 2 & 5 \\
\end{bmatrix}
=
\begin{bmatrix}
    1 & 0 \\
    2 & 1 \\
    3 & 2 \\
\end{bmatrix}
\begin{bmatrix}
    1 & 2 & 0 & 3 \\
    0 & 0 & 1 & -2 \\
\end{bmatrix}
$$
Interpretation: $\mathbf{v}_1 = 1\mathbf{c}_1 + 0\mathbf{c}_2,\dots, \mathbf{v}_4 = 3\mathbf{c}_1 + -2\mathbf{c}_2$. \\
Computation: Gauss-Jordan elimination (3.2.1). Get $R$ from RREF. $C$ is the columns from $A$ where there is a pivot in $R$.